\chapter{GIỚI THIỆU}
\label{chap:intro}

\section{Bài toán và động lực}

Các mô hình ngôn ngữ--thị giác (vision--language model, VLM) hiện đại như GPT-4V
hay LLaVA \autocite{liu2023llava} nhận đồng thời ảnh và văn bản, rồi sinh ra câu
trả lời bằng ngôn ngữ tự nhiên. Chúng đều được xây theo cùng một công thức:
lấy một mô hình ngôn ngữ lớn (LLM) đã tiền huấn luyện, ghép thêm một bộ mã hóa
thị giác, rồi tinh chỉnh để hai phương thức ``nói chuyện'' được với nhau. Công
thức này hiệu quả, nhưng với người học nó có một điểm mù: phần quan trọng
nhất --- \emph{cái LLM được tiền huấn luyện ra sao} --- luôn nằm trong hộp đen
của các tổ chức có hàng nghìn GPU.

Báo cáo này lấp điểm mù đó bằng cách làm ngược lại: \textbf{tự tiền huấn luyện
phần ngôn ngữ từ token đầu tiên}, ở quy mô mà một GPU 24\,GB xử lý được trong
khoảng nửa ngày, rồi mới ghép thị giác và đi tiếp các bước tinh chỉnh có giám
sát (SFT) và học tăng cường (RL). Sản phẩm là \textbf{ViVLM-nano}: một VLM
tiếng Việt 101 triệu tham số ngôn ngữ, nhận ảnh kèm câu lệnh tiếng Việt và sinh
mô tả hoặc câu trả lời tiếng Việt. Mọi tham số ngôn ngữ trong mô hình --- từ
embedding đến tầng attention cuối --- đều do chính pipeline trong báo cáo này
huấn luyện; thành phần huấn luyện sẵn duy nhất là bộ mã hóa thị giác SigLIP
\autocite{zhai2023siglip} được đóng băng, tương tự cách đồ án cuối kỳ dùng
ResNet đóng băng làm bộ trích đặc trưng.

Chọn tiếng Việt thay vì tiếng Anh vừa là thử thách vừa là đóng góp: dữ liệu web
tiếng Việt kém sạch hơn, chưa có bộ tách từ chuẩn ``mặc định'' như GPT-2, và các
VLM tiếng Việt hiện có (LaVy \autocite{tran2024lavy}, Vintern-1B
\autocite{doan2024vintern}) đều dựa trên LLM đa ngôn ngữ tỷ tham số huấn luyện
sẵn --- theo khảo sát của tôi, chưa có công bố nào đi trọn đường
\emph{tự pretrain tiếng Việt} $\rightarrow$ \emph{ghép thị giác} $\rightarrow$
\emph{SFT} $\rightarrow$ \emph{RL} ở quy mô một GPU.

\section{Vị trí đề tài trong sơ đồ kiến trúc đa phương thức của môn học}
\label{sec:intro-diagram}

Sơ đồ kiến trúc đa phương thức tổng quát của môn học gồm: các bộ mã hóa riêng
cho từng phương thức (image encoder, text encoder) sinh ra chuỗi biểu diễn
token, một khối \emph{hợp nhất} (fusion) bằng attention, đầu ra rẽ hai nhánh
\emph{phân loại} hoặc \emph{sinh chuỗi}, và --- phần được nhấn mạnh riêng cho
nhánh sinh --- chuỗi huấn luyện \textbf{SFT rồi RL}.

Ba bài làm của tôi phủ sơ đồ này theo ba bước tăng dần. Đồ án giữa kỳ
(Medical VQA) đi nhánh \emph{phân loại}: ảnh và câu hỏi được mã hóa, hợp nhất,
phân loại ra đáp án. Đồ án cuối kỳ (sinh mô tả ảnh Flickr8k) đi nhánh
\emph{sinh}, nhưng chỉ có một phương thức điều kiện (ảnh) và dừng ở teacher
forcing cộng một bước SCST. Báo cáo này lấp phần lõi còn lại: \emph{hai} luồng
token thị giác và ngôn ngữ cùng đổ vào một khối transformer duy nhất --- token
ảnh được chiếu vào không gian embedding của mô hình ngôn ngữ rồi ghép trước
chuỗi token văn bản, để self-attention tự làm nhiệm vụ hợp nhất --- và chuỗi
huấn luyện ba giai đoạn \emph{pretrain $\rightarrow$ SFT $\rightarrow$ RL} đúng
như vòng đời của các mô hình sinh hiện đại.

\section{Mục tiêu và câu hỏi nghiên cứu}

Đề tài đặt bốn câu hỏi, tương ứng bốn cụm thí nghiệm ở chương~\ref{chap:experiments}:
\begin{enumerate}
	\item Tiền huấn luyện từ đầu một mô hình 101M tham số kiến trúc Llama-style
	      trên khoảng 3 tỷ token tiếng Việt, ở quy mô một GPU, đạt được chất
	      lượng ngôn ngữ đến đâu (đường cong loss, perplexity, bits-per-character
	      so với baseline GPT-2 tiếng Việt huấn luyện sẵn)?
	\item Bộ tách từ BPE tự huấn luyện riêng cho tiếng Việt hiệu quả ra sao so
	      với tokenizer của PhoGPT và của GPT-2 tiếng Anh (đo bằng fertility ---
	      số token trung bình trên mỗi từ)?
	\item Ghép bộ mã hóa thị giác đóng băng vào mô hình ngôn ngữ tự pretrain qua
	      một projector pixel-shuffle, rồi SFT hai pha trên dữ liệu mô tả ảnh và
	      hỏi--đáp tiếng Việt, cho chất lượng đa phương thức đến đâu (CIDEr,
	      BLEU-4 trên UIT-ViIC, KTVIC, OpenViVQA)?
	\item SCST --- vốn được kiểm chứng ở đồ án cuối kỳ trên mô hình captioning
	      chuyên dụng --- còn phát huy tác dụng khi áp lên một VLM mà phần ngôn
	      ngữ tự pretrain không?
\end{enumerate}

Ràng buộc tự đặt, kế thừa từ hai đồ án trước: \emph{tự cài đặt} toàn bộ phần
học được --- khối transformer (RoPE, RMSNorm, SwiGLU, attention nhân quả),
vòng lặp huấn luyện ba giai đoạn, bộ tách từ, projector --- không dùng
\texttt{transformers} của Hugging Face cho phần ngôn ngữ (thư viện này chỉ được
dùng để nạp bộ mã hóa SigLIP đóng băng và mô hình baseline khi đánh giá).

\section{Đóng góp và cấu trúc báo cáo}

Đóng góp chính của báo cáo: (i) một pipeline hoàn chỉnh, kiểm thử được
(52 unit test chạy ngoại tuyến), tái lập được, đưa một mô hình ngôn ngữ tiếng
Việt từ dữ liệu web thô đến VLM biết trả lời bằng tiếng Việt ở quy mô một GPU;
(ii) các đo đạc định lượng về chất lượng ngôn ngữ, hiệu quả tokenizer và chất
lượng đa phương thức ở quy mô ``nano''; (iii) phân tích mối liên hệ giữa bộ
tách từ bản ngữ và hiệu quả nén (bits-per-character) --- cho thấy vì sao một
mô hình nhỏ tự pretrain vẫn vượt baseline về nén ký tự tiếng Việt.

Báo cáo gồm 5 chương. Chương~\ref{chap:background} trình bày cơ sở lý thuyết:
mô hình ngôn ngữ tự hồi quy, các thành phần kiến trúc Llama-style, BPE, quy
luật scaling, hợp nhất thị giác kiểu LLaVA, SFT và SCST, cùng các công trình
liên quan. Chương~\ref{chap:methods} mô tả dữ liệu, kiến trúc và cấu hình
huấn luyện ba giai đoạn. Chương~\ref{chap:experiments} trình bày kết quả định
lượng và định tính kèm bàn luận. Chương~\ref{chap:conclusion} kết luận, nêu
hạn chế và hướng phát triển.
Chương~\ref{chap:conclusion} kết luận, nêu hạn chế và hướng phát triển.
